\section{Full-Wave Numerical Verification}

\subsection{Convergence and resolution checks}

Within the declared local annulus and on the tracked branch, the isotropic
coordinates and curvature passed the predeclared order-convergence gates.
The eigenpairs and continuation history separately passed the residual,
orthogonality, overlap, and eigengap gates.  The independently computed values
remained consistent with
\(\kappa_0=\ZGVKappa\), \(\Omega_0=\ZGVOmega\), and
\(a=\ZGVCurvature\), while Fig.~\ref{fig:isotropic-zgv} exposes the full
order sequence rather than only its finest member.  This evidence concerns the
selected local branch and does not assert completeness for the dispersion
spectrum outside the configured window.

No coordinate match was accepted alone.  A converged row also had to satisfy
the normalized residual of Eq.~\eqref{eq:normalized-eigen-residual}, the mass
orthogonality defect, the mass-weighted overlap of Eq.~\eqref{eq:mass-mac}, and
the relative separation gate in Eq.~\eqref{eq:eigengap-rejection}.  The
generated Supplementary convergence table reports the final one- and
two-element values beside their declared gates.  Algorithm~\ref{alg:mode-tracking}
supplied the branch history, so apparent coordinate agreement after a branch
switch could not pass this joint check \citep{adamou2004spectral,kiefer2023computing}.

The anisotropic evidence closed a distinct joint gate.  Centered
full-anisotropic differences checked \(V_4=\VFour\) and the differentiated-mode
equation, while Hessian discrepancy estimates and candidate-grid refinement
checked the \(\MorseMinimumCount\)-minimum and
\(\MorseSaddleCount\)-saddle set in
Figs.~\ref{fig:angular-sensitivity} and \ref{fig:morse-points}.  The two
resolution routes are defined in Eqs.~\eqref{eq:differentiated-mode-solve}
and \eqref{eq:critical-point-uncertainty}.  Boundary noncriticality and the
index identity of Eq.~\eqref{eq:annular-index-certificate} then closed the
declared local annulus.  Together these checks support a resolved local set
on one branch, not a global enumeration or laboratory evidence.

\subsection{Response and phase accuracy}

For the declared joint-limit rows and the separate fixed-anisotropy row,
independent full-wave polar-grid refinements converged while the maximum
accumulated inter-resolution phase-discrepancy estimator remained
\(\MaxPhaseError\), below its declared acceptance threshold.  Each production
grid rebuilt and retracked its own spectral surface
before applying the Cartesian-measure quadrature in
Eq.~\eqref{eq:polar-green-quadrature}.  The nested subsampling stored in the
auxiliary convergence artifact is therefore a diagnostic only, not a
substitute for these independent production surfaces.

The accumulated inter-resolution phase-discrepancy estimator in
Eq.~\eqref{eq:phase-error-certificate} combines the local nested-order
frequency difference with the common-node difference between successive polar
grids.  It restricts the accepted reported window by comparing computed
resolutions; without a continuum enclosure, it does not bound the unknown
continuum phase error.  Figure~\ref{fig:decay-crossover} shows that the complex
response and fixed-anisotropy comparison passed this same declared consistency
gate.

The masks fixed before response evaluation in Eq.~\eqref{eq:fit-masks} gave an
early decay slope
\(\EarlyDecaySlope\) and a late decay slope \(\LateDecaySlope\).  The early
fit uses only the smallest perturbation row before angular stationary points
are resolved, whereas the late fit uses complete half-open root-mean-square
modulation bins in the fixed-anisotropy row.  The fitted log--log intercept is
a nuisance parameter for exponent estimation; it is not an adjustment of the
transition or exact-Morse response curves, which use no fitted alignment.

\subsection{Source and window robustness}

The declared robustness sweep varied the Gaussian source radius only within
its declared range.  A separate sweep varied the width of the local window in
wavenumber.  The response norm changed by a bounded amount, while the baseline
normalization and localization at the widest window boundary were preserved.
Every declared pair was evaluated from the stored branch surface.  This check
concerns source and truncation sensitivity, rather than a new decay exponent or
critical-point completeness.  The full sweep and its boundary-weight
diagnostic are reported in the Supplementary Information.

The baseline row is included in both one-parameter sweeps, so every reported
difference has a common normalization.  The window check also bounds the
weighted contribution at the radial endpoints, limiting the possibility that
the response trend is created by abrupt branch truncation.  These controls
support the interpretation of Fig.~\ref{fig:decay-crossover} without replacing
the asymptotic hypotheses or the independent phase-discrepancy check.

\subsection{Silicon stress test}

A separate calculation for a \([0\,0\,1]\)-cut silicon plate recovered an
alternating local minimum--saddle set with resolved Cartesian Hessians and noncritical
annular boundaries, serving only as a finite-anisotropy stress test.  The
calculation repeats the tracked-branch search and the winding procedure of
Eq.~\eqref{eq:annular-index-certificate}, but physical silicon lies outside the
controlled weak-anisotropy sequence used for the coefficient and remainder
claims \citep{prada2009anisotropy,kiefer2023beating}.

Accordingly, this result is reported as numerical robustness of the search,
classification, and boundary-check machinery under a materially larger
constitutive contrast.  It supplies neither evidence for the weak-parameter
scaling laws nor a count over the complete silicon dispersion surface.  The
local alternating organization is consistent with earlier anisotropic ZGV
computations, while the claims of this study remain attached to the declared
perturbative family and its explicitly tracked branch.
The corresponding contours, refined points, Hessian inertia, and boundary
margin are reported in the Supplementary Information.
